\documentclass[10pt,a4paper]{amsart}
\usepackage{amsmath,amssymb,amsthm,graphicx}
% Retained minimal preamble for the Stacks Project chapter varieties.tex.
% Verbatim excerpt of the chapter preamble: Xy-pic setup (native preamble lines 14-18), the native theorem
% declarations with the mathematical macros (native preamble lines 154-217), and the project own no-op
% definitions of the reference/slogan/history display environments (native wrapper.tex lines 1-3) that the
% counted statement uses for its slogan. Copied mechanically; nothing added, removed or redefined.
% For commutative diagrams we use Xy-pic
\usepackage[all]{xy}

% We use 2cell for 2-commutative diagrams.
\xyoption{2cell}
% Theorem environments.
%
\theoremstyle{plain}
\newtheorem{theorem}[subsection]{Theorem}
\newtheorem{proposition}[subsection]{Proposition}
\newtheorem{lemma}[subsection]{Lemma}

\theoremstyle{definition}
\newtheorem{definition}[subsection]{Definition}
\newtheorem{example}[subsection]{Example}
\newtheorem{exercise}[subsection]{Exercise}
\newtheorem{situation}[subsection]{Situation}

\theoremstyle{remark}
\newtheorem{remark}[subsection]{Remark}
\newtheorem{remarks}[subsection]{Remarks}

\numberwithin{equation}{subsection}

% Macros
%
\def\lim{\mathop{\mathrm{lim}}\nolimits}
\def\colim{\mathop{\mathrm{colim}}\nolimits}
\def\Spec{\mathop{\mathrm{Spec}}}
\def\Hom{\mathop{\mathrm{Hom}}\nolimits}
\def\Ext{\mathop{\mathrm{Ext}}\nolimits}
\def\SheafHom{\mathop{\mathcal{H}\!\mathit{om}}\nolimits}
\def\SheafExt{\mathop{\mathcal{E}\!\mathit{xt}}\nolimits}
\def\Sch{\mathit{Sch}}
\def\Mor{\mathop{\mathrm{Mor}}\nolimits}
\def\Ob{\mathop{\mathrm{Ob}}\nolimits}
\def\Sh{\mathop{\mathit{Sh}}\nolimits}
\def\NL{\mathop{N\!L}\nolimits}
\def\CH{\mathop{\mathrm{CH}}\nolimits}
\def\proetale{{pro\text{-}\acute{e}tale}}
\def\etale{{\acute{e}tale}}
\def\QCoh{\mathit{QCoh}}
\def\Ker{\mathop{\mathrm{Ker}}}
\def\Im{\mathop{\mathrm{Im}}}
\def\Coker{\mathop{\mathrm{Coker}}}
\def\Coim{\mathop{\mathrm{Coim}}}

% Boxtimes
%
\DeclareMathSymbol{\boxtimes}{\mathbin}{AMSa}{"02}

%
% Macros for moduli stacks/spaces
%
\def\QCohstack{\mathcal{QC}\!\mathit{oh}}
\def\Cohstack{\mathcal{C}\!\mathit{oh}}
\def\Spacesstack{\mathcal{S}\!\mathit{paces}}
\def\Quotfunctor{\mathrm{Quot}}
\def\Hilbfunctor{\mathrm{Hilb}}
\def\Curvesstack{\mathcal{C}\!\mathit{urves}}
\def\Polarizedstack{\mathcal{P}\!\mathit{olarized}}
\def\Complexesstack{\mathcal{C}\!\mathit{omplexes}}
% \Pic is the operator that assigns to X its picard group, usage \Pic(X)
% \Picardstack_{X/B} denotes the Picard stack of X over B
% \Picardfunctor_{X/B} denotes the Picard functor of X over B
\def\Pic{\mathop{\mathrm{Pic}}\nolimits}
\def\Picardstack{\mathcal{P}\!\mathit{ic}}
\def\Picardfunctor{\mathrm{Pic}}
\def\Deformationcategory{\mathcal{D}\!\mathit{ef}}
\newenvironment{reference}{}{}
\newenvironment{slogan}{}{}
\newenvironment{history}{}{}
\usepackage{cleveref}
\crefname{originallemma}{original lemma}{original lemmas}
\Crefname{originallemma}{Original lemma}{Original lemmas}
\makeatletter
\newlabel{algebra-lemma-Noether-normalization}{{00OY}{source}{}{sourcepointer-1}{}}
\newlabel{algebra-lemma-Noether-normalization@cref}{{[originallemma][0][]00OY}{[1][0][]source}{}{}{}}
\makeatother
\crefname{originallemma}{original lemma}{original lemmas}
\Crefname{originallemma}{Original lemma}{Original lemmas}
\makeatletter
\newlabel{algebra-lemma-artinian-finite-length}{{00JB}{source}{}{sourcepointer-2}{}}
\newlabel{algebra-lemma-artinian-finite-length@cref}{{[originallemma][0][]00JB}{[1][0][]source}{}{}{}}
\makeatother
\crefname{originallemma}{original lemma}{original lemmas}
\Crefname{originallemma}{Original lemma}{Original lemmas}
\makeatletter
\newlabel{algebra-lemma-finite-dimensional-algebra}{{00J6}{source}{}{sourcepointer-3}{}}
\newlabel{algebra-lemma-finite-dimensional-algebra@cref}{{[originallemma][0][]00J6}{[1][0][]source}{}{}{}}
\makeatother
\crefname{originallemma}{original lemma}{original lemmas}
\Crefname{originallemma}{Original lemma}{Original lemmas}
\makeatletter
\newlabel{schemes-lemma-disjoint-union-affines}{{01I5}{source}{}{sourcepointer-4}{}}
\newlabel{schemes-lemma-disjoint-union-affines@cref}{{[originallemma][0][]01I5}{[1][0][]source}{}{}{}}
\makeatother
\title{Zero dimensional locally algebraic schemes are disjoint unions of spectra of local Artinian algebras}
\author{The Stacks Project Authors}
\date{stacks-project:varieties.tex:local-snapshot-20261011; GNU FDL 1.2 or later}
\begin{document}
\maketitle
\noindent\textit{Editorial note.} All delivered mathematical text is the author's own, copied from the pinned Stacks Project chapter file varieties.tex; no statement or proof step was rewritten, reordered or removed.\par\medskip
\noindent\textit{Editorial note.} The proof is complete as written and ends with the finiteness of X over Spec(k); no step is referred to the reader.\par\medskip
\noindent\textit{Editorial note.} The counted spans carry 4 navigation stubs and no bibliography citation.\par\medskip
\noindent\textit{Editorial note.} The preamble is the chapter's retained minimal preamble (sources/stacks-varieties/retained-preamble.tex lines 1-77), a verbatim excerpt of the native preamble.\par\medskip
\noindent\textit{Editorial note.} The label lemma-algebraic-scheme-dim-0 is defined only in varieties.tex, so no cross-chapter restatement had to be weighed for this card.\par\medskip
\noindent\textit{Allowed standard prerequisites (editorial).}\begin{itemize}
\item Noether normalization for a zero dimensional algebra over a field. Algebra, Lemma algebra-lemma-Noether-normalization, invoked at proof lines 3748-3750. Supplies the finite injection k to A, hence finite k-dimension of A.
\item A finite dimensional algebra over a field is Artinian and a product of Artinian local rings. Algebra, Lemmas algebra-lemma-finite-dimensional-algebra and algebra-lemma-artinian-finite-length, invoked at proof lines 3751-3755. Supplies the product decomposition of A into Artinian local rings.
\item A finite disjoint union of affine schemes is affine. Schemes, Lemma schemes-lemma-disjoint-union-affines, invoked at proof lines 3772-3773. Supplies affineness of a finite disjoint union of spectra in the algebraic case.
\end{itemize}\medskip
\ifdefined\hypertarget\hypertarget{sourcepointer-1}{}\fi
\noindent\textit{Source reference.} Original lemma 00OY: sources/stacks-chow/tags.txt, source lines 912--912. Pointer outside the counted span, invoked once for context; the load-bearing argument of the counted statement is the author's proof reproduced in full in the proof span.\par\medskip
\ifdefined\hypertarget\hypertarget{sourcepointer-2}{}\fi
\noindent\textit{Source reference.} Original lemma 00JB: sources/stacks-chow/tags.txt, source lines 709--709. Pointer outside the counted span, invoked once for context; the load-bearing argument of the counted statement is the author's proof reproduced in full in the proof span.\par\medskip
\ifdefined\hypertarget\hypertarget{sourcepointer-3}{}\fi
\noindent\textit{Source reference.} Original lemma 00J6: sources/stacks-chow/tags.txt, source lines 704--704. Pointer outside the counted span, invoked once for context; the load-bearing argument of the counted statement is the author's proof reproduced in full in the proof span.\par\medskip
\ifdefined\hypertarget\hypertarget{sourcepointer-4}{}\fi
\noindent\textit{Source reference.} Original lemma 01I5: sources/stacks-chow/tags.txt, source lines 1966--1966. Pointer outside the counted span, invoked once for context; the load-bearing argument of the counted statement is the author's proof reproduced in full in the proof span.\par\medskip
\begin{lemma}
\label{lemma-algebraic-scheme-dim-0}
Let $k$ be a field. Let $X$ be a locally algebraic $k$-scheme of
dimension $0$. Then $X$ is a disjoint union of spectra of local Artinian
$k$-algebras $A$ with $\dim_k(A) < \infty$. If $X$ is an algebraic $k$-scheme
of dimension $0$, then in addition $X$ is affine and the morphism
$X \to \Spec(k)$ is finite.
\end{lemma}

\begin{proof}
Let $X$ be a locally algebraic $k$-scheme of dimension $0$.
Let $U = \Spec(A) \subset X$ be an affine open subscheme.
Since $\dim(X) = 0$ we see that $\dim(A) = 0$.
By Noether normalization, see
Algebra, Lemma \ref{algebra-lemma-Noether-normalization}
we see that there exists a finite injection $k \to A$, i.e.,
$\dim_k(A) < \infty$. Hence $A$ is Artinian, see
Algebra, Lemma \ref{algebra-lemma-finite-dimensional-algebra}.
This implies that $A = A_1 \times \ldots \times A_r$ is a product
of finitely many Artinian local rings, see
Algebra, Lemma \ref{algebra-lemma-artinian-finite-length}.
Of course $\dim_k(A_i) < \infty$ for each $i$ as the sum of
these dimensions equals $\dim_k(A)$.

\medskip\noindent
The arguments above show that $X$ has an open covering whose members are
finite discrete topological spaces. Hence $X$ is a discrete topological space.
It follows that $X$ is isomorphic to the disjoint union of its connected
components each of which is a singleton. Since a singleton scheme is affine
we conclude (by the results of the paragraph above) that each of these
singletons is the spectrum of a local Artinian $k$-algebra $A$ with
$\dim_k(A) < \infty$.

\medskip\noindent
Finally, if $X$ is an algebraic $k$-scheme of dimension $0$, then
$X$ is quasi-compact hence is a finite disjoint union
$X = \Spec(A_1) \amalg \ldots \amalg \Spec(A_r)$
hence affine (see
Schemes, Lemma \ref{schemes-lemma-disjoint-union-affines})
and we have seen the finiteness of $X \to \Spec(k)$ in the
first paragraph of the proof.
\end{proof}

\end{document}
